% Chapter 2, Topic _Linear Algebra_ Jim Hefferon
%  http://joshua.smcvt.edu/linearalgebra
%  2001-Jun-11
\topic{域}
\index{field|(}
只涉及整数或只涉及有理数的计算，
比涉及实数的计算容易得多。
其他代数结构，比如整数或有理数，
能否在向量空间的定义中
取代 \( \Re \) 的位置？

如果把“能起作用”理解为本章的结论
仍然成立，
那么就有一份自然的条件清单，
一个结构（即数系）
必须满足这些条件才能取代 $\Re$ 的位置。
\definend{域}\index{field!definition} 是一个集合
\( \F \)，连同运算
`\( + \)'
与`\( \cdot \)'，满足下列条件：
\begin{tfae}
   \item 对任意 \( a,b\in\F \)，\( a+b \) 的结果仍在 \( \F \) 中，并且
         \( a+b=b+a \)，并且
         若 \( c\in\F \)，则 \( a+(b+c)=(a+b)+c \)
   \item 对任意 \( a,b\in\F \)，\( a\cdot b \) 的结果仍在 \( \F \) 中，并且
         \( a\cdot b=b\cdot a \)，并且
         若 \( c\in\F \)，则 \( a\cdot (b\cdot c)=(a\cdot b)\cdot c \)
   \item 若 \( a,b,c\in\F \)，则 \( a\cdot (b+c)=a\cdot b+a\cdot c \)
   \item 存在元素 \( 0\in\F \)，使得
         若 \( a\in\F \) 则 \( a+0=a \)，并且
         对每个 \( a\in\F \)，存在元素 \( -a\in\F \)
                使得 \( (-a)+a=0 \)
   \item 存在元素 \( 1\in\F \)，使得
         若 \( a\in\F \) 则 \( a\cdot 1=a \)，并且
         对 \( \F \) 中每个元素 \( a\neq 0 \)，
                存在元素 \( a^{-1}\in\F \)
                使得 \( a^{-1}\cdot a=1 \)。
\end{tfae}

例如，由实数集连同其通常的
加法与乘法运算构成的代数结构就是一个域。
另一个域是带有其通常加法
与乘法运算的有理数集。
不是域的代数结构的例子是
整数，因为它不满足最后一个条件。

有些例子更出人意料。
集合 \( \mathbb{B}=\set{0,1} \) 配上以下运算：
\begin{center}
  \begin{tabular}{c|cc}
    \( + \) &\( 0 \) &\( 1 \) \\
    \hline
    \( 0 \) &\( 0 \) &\( 1 \) \\
    \( 1 \) &\( 1 \) &\( 0 \)
  \end{tabular}
  \qquad
  \begin{tabular}{c|cc}
    \( \cdot \) &\( 0 \) &\( 1 \) \\
     \hline
       \( 0 \)  &\( 0 \) &\( 0 \)  \\
       \( 1 \)  &\( 0 \) &\( 1 \)
  \end{tabular}
\end{center}
是一个域；见\nearbyexercise{exer:BinField}。

在本书中，我们本可以把线性代数发展成
标量取自任意域的向量空间的理论。
在那种情况下，
把`\( \Re \)'换成`\( \F \)'，也就是
把系数、向量分量
和矩阵元素都取为 \( \F \) 中的元素，那么这里几乎所有的命题都会照搬过来
（例外是涉及距离或夹角的命题，
它们需要另外的展开）。
下面是一些例子；每一条都对域 \( \F \) 上的向量空间 \( V \) 成立。
\begin{itemize}
  %\makebox[\parindent]{\ \hfil\ }
  \item[$*$] 对任意 \( \vec{v}\in V \) 与 \( a\in\F \)，
  (i)~\( 0\cdot\vec{v}=\zero \)，
  (ii)~\( -1\cdot\vec{v}+\vec{v}=\zero \)，
  以及 (iii)~\( a\cdot\vec{0}=\vec{0} \)。

  \item[$*$] \( V \) 的一个子集的张成，即线性组合的集合，
  是 \( V \) 的子空间。

  \item[$*$] 线性无关集的任何子集也
     线性无关。

  \item[$*$] 在有限维向量空间中，
    任意两个基所含元素的个数相同。
\end{itemize}
（即使是没有明确提到 \( \F \) 的命题，
其证明中也用到了域的性质。）

我们不会在这一更一般的框架下发展向量空间理论，因为
额外的抽象可能使人分心。
我们想要阐明的那些想法，在我们坚持使用实数时就已经出现了。

例外是第五章。
在那里我们必须对多项式作因式分解，
所以我们将转而考虑
复数域上的向量空间。

\begin{exercises}
  \item 
    验证实数构成一个域。
    \begin{answer}
      逐条检查这五个条件可知，它们都是初等数学中我们熟悉的内容。
    \end{answer}
  \item 
    证明下述结构都是域。
    \begin{exparts*}
       \partsitem 有理数 $\Q$
       \partsitem 复数 $\C$
    \end{exparts*}
    \begin{answer}
      与前面的问题一样，逐条检查这五个条件
      可知，
      对这两个结构而言，
      这些性质都是我们熟悉的。
    \end{answer}
  \item 
     举一个例子，说明整数数系
     不是域。
     \begin{answer}
       整数不满足条件~(5)。
       例如，$2$ 没有乘法逆元\Dash 虽然
       $2$ 是整数，$1/2$ 却不是。
     \end{answer}
  \item \label{exer:BinField} 
     验证集合 $\mathbb{B}=\set{0,1}$ 在上面所列的运算下是一个域，
     \begin{answer}
       我们可以通过列举所有可能的情况来做这些验证。
       例如，为验证条件~(2)的前半部分，我们必须检查
       该结构对加法封闭，并且加法
       满足交换律 $a+b=b+a$，
       由于这样的数对
       只有四对，我们可以对所有可能的数对 $a$ 与~$b$ 检验这两点。
       类似地，对于结合律，三元组 $a$、$b$、
       $c$ 只有八种，因此检验并不太长。
       （还有其他做这些验证的方法；特别地，你可能会
       认出这些运算就是模~$2$ 的算术。
       但穷举检验也并不繁重）
     \end{answer}
  \item 
     给集合 $\set{0,1,2}$
     配上适当的运算使其成为一个域。     
     \begin{answer}
       下面这些运算就可以。
       \begin{center}
          \begin{tabular}{c|ccc}
             \( + \) &\( 0 \) &\( 1 \) &$2$ \\
             \hline
             \( 0 \) &\( 0 \) &\( 1 \) &$2$ \\
             \( 1 \) &\( 1 \) &\( 2 \) &$0$ \\
             $2$     &$2$     &$0$     &$1$
          \end{tabular}
          \qquad
          \begin{tabular}{c|ccc}
            \( \cdot \) &\( 0 \) &\( 1 \) &$2$ \\
            \hline
            \( 0 \)  &\( 0 \) &\( 0 \) &$0$ \\
            \( 1 \)  &\( 0 \) &\( 1 \) &$2$ \\
            $2$      &$0$     &$2$     &$1$
          \end{tabular}
       \end{center}
       与前面的题目一样，我们可以
       通过列举所有情况来
       验证它们满足这些条件。
     \end{answer}
\end{exercises}
\index{field|)}
%\include{la}
%\end{document}
